\section{Auxiliary decision-theoretic bounds}
\label{app:decision-bounds}

This appendix collects finite-sample decision bounds and illustrative
optimization examples. They support the interpretation of the assembly
results but are not additional structural hypotheses.

\begin{remark}[Finite-sample consistency and the scope of an error bound]
\label{rem:sample-complexity}
For a truth $S$ and finite competitor family $\mathcal C_S$ with iid
read laws $p_D$ distinct from $p_S$, let
$\rho(p,q)=\sum_w\sqrt{p(w)q(w)}$. The Hellinger union bound gives
\[
 \Pr_S\{\exists D\in\mathcal C_S:L_D(R_1,\ldots,R_n)
                                  \ge L_S(R_1,\ldots,R_n)\}
 \le \sum_{D\in\mathcal C_S}\rho(p_S,p_D)^n
 \le |\mathcal C_S|\,\rho_*^{\,n}.
\]
Here $\rho_*<1$ when the competitor family is nonempty and identifiable.
It proves an exponential \emph{rate}, not necessarily a practical sample
size: the class may be enormous and $\rho_*$ close to one.

For arbitrary same-length circular genomes the uniform strict-separation
assumption fails: at $L=2$, $AAABAB$ and $AABAAB$ have identical read laws
(Example~\ref{ex:ABAB}), giving $\rho=1$ for a family containing both.
The boundary is sharp and kernel-checked: when the competitor family
contains the truth itself, every read tuple ties for the exact-ML lead, so
the exact wrong-selection probability is $1$ at every sample size $n$
(\path{HellingerNoDecay.eventProb_wrong_eq_one_of_mem}); no decay rate is
possible without identifiability.
In contrast, Theorem~\ref{thm:p2-finite-class} proves that the entire
primitive-P2 candidate class is finite for fixed alphabet and read length,
and Theorem~\ref{thm:population} proves separation modulo rotation.
Consequently Corollary~\ref{cor:p2-finite-recovery-all} applies the bound
to that complete class, not merely to an arbitrarily truncated shortlist.
Numerical values of its worst-case affinity are not computed here.
See Theorem~\ref{thm:finite-recovery} for the kernel-checked proof.
\end{remark}

\begin{corollary}[Explicit sufficient sample size]
\label{cor:explicit-sample-size}
In the setting of Remark~\ref{rem:sample-complexity}, suppose
$\mathcal C_S$ is finite and nonempty with uniform affinity
$\rho_*=\max_{D\in\mathcal C_S}\rho(p_S,p_D)<1$. Then for every
$\delta\in(0,1]$,
\[
  n\;\ge\;
  \left\lceil\frac{\ln(|\mathcal C_S|/\delta)}{\ln(1/\rho_*)}\right\rceil
  \quad\Longrightarrow\quad
  \Pr_S\!\left\{\exists D\in\mathcal C_S:
      L_D(R_1,\dots,R_n)\ge L_S(R_1,\dots,R_n)\right\}
  \;\le\;\delta .
\]
For ML selection over $\{S\}\cup\mathcal C_S$, the truth's rotation class is
then the strict maximizer with probability at least $1-\delta$.
\end{corollary}
\begin{proof}
The choice of $n$ gives
$|\mathcal C_S|\,\rho_*^{\,n}
=|\mathcal C_S|\,e^{-n\ln(1/\rho_*)}\le\delta$, which inserted into the union
bound of Remark~\ref{rem:sample-complexity} proves the display.
\end{proof}
The threshold is a \emph{sufficient}, conservative sample size: it inverts a
union bound, is not necessary or minimax-optimal, and inherits the
finite-family and uniform-separation hypotheses. The Hellinger constant is the
one kernel-checked in \texttt{AssemblyP1.FiniteReadRecovery}, whose concrete
instance \texttt{il\_confidence\_99} is kernel-checked there. The next
proposition sharpens the rate from the Hellinger constant to the Chernoff
exponent, and the remark after it records how far the resulting bound remains
from the exact finite-sample error.

\begin{proposition}[One-sided likelihood-ratio refinement]
\label{prop:chernoff-refinement}
In the setting of Remark~\ref{rem:sample-complexity}, let
$\mathcal C_S$ be a finite nonempty family of competitors whose read laws
$d=p_D$ differ from $p=p_S$. For any distributions $p,q$ on the
finite read alphabet $W$, put
\[
 M_{p,q}(t)=\sum_{w:\,p(w)>0}
     p(w)\left(\frac{q(w)}{p(w)}\right)^t,\qquad t>0,
 \qquad
 E_+(p,q)=-\ln\inf_{t>0}M_{p,q}(t),
\]
where $-\ln 0=+\infty$ and $0^t=0$ for $t>0$.
Set $E_*=\min_{D\in\mathcal C_S}E_+(p_S,p_D)$.
Then $E_+(p,q)\ge-\ln\rho(p,q)$, so
$E_*\ge\ln(1/\rho_*)>0$; and, for every integer $n\ge1$,
\[
 \Pr_S\{\exists D\in\mathcal C_S:
          L_D(R_1,\ldots,R_n)\ge L_S(R_1,\ldots,R_n)\}
 \;\le\;\sum_{D\in\mathcal C_S}e^{-nE_+(p_S,p_D)}
 \;\le\;|\mathcal C_S|e^{-nE_*}.
\]
For $E_*<\infty$ it therefore suffices to take
$n\ge\max\{1,\lceil\ln(|\mathcal C_S|/\delta)/E_*\rceil\}$
to obtain error at most $\delta\in(0,1]$.
If $E_*=\infty$, every competitor has zero chance of tying or
beating the truth for every $n\ge1$.

The exponent $E_+$ is \emph{one-sided}: the minimization ranges over
$t>0$, not just $t\in[0,1]$ as in the conventional symmetric
Chernoff information $C_{\rm Ch}(p,q)$.
If $\supp q\subseteq\supp p$, the minimizer may be taken in
$[0,1]$ (with boundary limits), so
$E_+(p,q)=C_{\rm Ch}(p,q)$. Without that support condition, the
two exponents can differ, possibly with $E_+=\infty$.
\end{proposition}
\begin{proof}
Under $p_S$, $L_S>0$ almost surely. For any $t>0$, independence gives
$\mathbb E_S[(L_D/L_S)^t]=M_{p_S,p_D}(t)^n$.
On the event $L_D\ge L_S$, the likelihood ratio raised to $t$ is at
least one. Markov's inequality, the infimum over $t>0$, and the union
bound give the display. The choice $t=1/2$ gives
$E_+(p,q)\ge-\ln\rho(p,q)$; hence $E_*>0$ by the strict
identifiability hypothesis. If $E_*=\infty$ every term in the
bound vanishes for $n\ge1$.

When $\supp q\subseteq\supp p$, the log moment function
$\ln M_{p,q}(t)$ is convex on $(0,\infty)$ and
its left derivative at $t=1$ equals
$\sum_{w:q(w)>0}q(w)\ln(q(w)/p(w))
=\KL(q\,\|\,p)\ge0$. Therefore its infimum over $t>0$
is achieved or approached with $t\le1$, giving the usual Chernoff
expression on this support-restricted domain.
\end{proof}
\statusline{Paper-level application of the classical one-sided
large-deviation (Cram\'er) bound \parencite{cover2006,tsybakov2009}.
The Hellinger-based theorem in
\texttt{AssemblyP1.FiniteReadRecovery} remains kernel-checked;
this refinement and the large-deviation sharpness assertion below
are not Lean theorems.}

\begin{remark}[The sharp one-sided exponent and a support-mismatch example]
\label{rem:chernoff-sharp}
For a finite read alphabet and the event that a \emph{wrong} candidate
ties or beats the truth, the exact pairwise large-deviation exponent is
\[
 \lim_{n\to\infty}-\frac1n\ln
   \Pr_p\{L_q(R_1,\ldots,R_n)\ge L_p(R_1,\ldots,R_n)\}
 = E_+(p,q),
\]
using $-\ln0=\infty$. This is the one-sided Cram\'er/Sanov rate
for the log-likelihood ratio, including the boundary cases where
the event is impossible. For a finite competitor family, the exponent
of its union event is $E_*=\min_D E_+(p_S,p_D)$:
the probability is between its largest pairwise probability and the
sum of those probabilities. The finite-$n$ union bound is generally
loose because the events overlap and polynomial factors are omitted.

The restriction $t\in[0,1]$ in the \emph{symmetric} Chernoff
information cannot in general be used for this one-sided claim.
Consider the circular genomes
$S=\texttt{AAB}$ and $D=\texttt{AABB}$ at read length $L=2$.
The possible oriented reads of $S$ are
$\texttt{AA},\texttt{AB},\texttt{BA}$, each with probability $1/3$.
Under $D$ those three reads each have probability $1/4$, and
the additional read $\texttt{BB}$ also has probability $1/4$.
For \emph{every} sample of $n\ge1$ reads from $S$,
\[
 \frac{L_D}{L_S}=(3/4)^n<1.
\]
Thus the exact wrong-candidate event has probability zero and exponent
$E_+(p_S,p_D)=\infty$, because $M_{p_S,p_D}(t)=(3/4)^t$ tends
to zero as $t\to\infty$. In contrast,
$C_{\rm Ch}(p_S,p_D)=\ln(4/3)$.
The source of the discrepancy is the read $\texttt{BB}$, which
$D$ can generate but $S$ cannot.

On pairs with $\supp q\subseteq\supp p$, the one-sided and
conventional Chernoff exponents agree. For example, in the
two-symbol worked family of
Appendix~\ref{app:reproduction}, the exact error
at $n=16$ is about $3.7$ times smaller than the union bound,
even though their exponential rates agree. The script
Appendix~\ref{app:reproduction}
also checks the support-mismatch witness above.
\end{remark}

\begin{remark}[Exact finite-$n$ sample size for one explicit family]
\label{rem:minimax-sample-size}
The union bounds above are sufficient, not tight. For the single explicit
family $A=(\tfrac12,\tfrac12)$, $B=(\tfrac34,\tfrac14)$, $C=(\tfrac14,\tfrac34)$
on the binary read alphabet, the exact finite-$n$ optima can be computed in
closed form. Under the uniform prior the Bayes-optimal (prior-averaged) success
$V_n=\sum_k\max_t\tfrac13 P_t(k)$ is non-decreasing and reaches $99\%$ exactly
at $n=85$ ($V_{84}<0.99\le V_{85}$); the minimax-over-truths error
$\min_d\max_t P_t(d\text{ wrong})$, equal to the maximum Bayes error over
priors, falls to $1\%$ exactly at $n=87$. Both are tight: the Bayes risk is the
minimum possible prior-averaged error, and the minimax value is certified by
the least favourable symmetric prior together with the optimally tie-split
Bayes rule. The per-truth success at $n=85$ is $0.98349\ldots$ for the middle
truth $A$ and $0.99336\ldots$ for $B,C$, so the $99\%$ figure is prior-averaged,
not per-truth. The kernel-checked Hellinger union bound instead gives the
sufficient size $153$, about $1.8\times$ the exact optimum. These numbers are
exact finite arithmetic for this one family, not a Lean result and not a
uniform sample size over the genome universe; the decision-theoretic
ingredients are standard \parencite{berger1985,lehmann1998}. See
Appendix~\ref{app:reproduction}, checked by
Appendix~\ref{app:reproduction}.
\end{remark}

\begin{corollary}[Conditional error while bridging holds]
\label{cor:conditional-bridging-error}
Let $B_n$ be the event that the true genome satisfies a specified bridging
predicate on $n$ iid reads, with $\Pr_S(B_n)>0$. For any finite candidate
family $\mathcal C_S$ as in Remark~\ref{rem:sample-complexity},
\[
 \Pr_S(\text{a wrong candidate ties or beats }S\mid B_n)
 \le
 \frac{\sum_{D\in\mathcal C_S}\rho(p_S,p_D)^n}
      {\Pr_S(B_n)}.
\]
\end{corollary}
\begin{proof}
For any event $E$, $\Pr(E\mid B_n)\le\Pr(E)/\Pr(B_n)$.
Apply this to the wrong-candidate event and use the preceding union bound.
\end{proof}
The conditioning here is on the \emph{truth's} bridging event. It is not
a switch to Question~A's restricted posterior, and it becomes informative
only after bounding the denominator and the pairwise affinities.
